% =============================================================================
% 03-system-calls-notes.tex — Lecture Notes for Week 3:
% System Calls — Mechanics, Parameter Passing, and Types
%
% DERIVED ARTIFACT. Source of truth: Slides/CS502/03-system-calls.tex.
% Generated by /lecture-notes; do not hand-edit content independently of the
% Beamer source without re-running /qa-notes afterward (single-source-of-truth.md).
% Structured as a textbook chapter (numbered sections/figures/examples,
% topic-organized rather than a frame-by-frame transcript).
%
% Anchor reading: Silberschatz, Galvin & Gagne, Operating System Concepts,
%                 8th ed. — Ch. 2 §2.3 (system calls, pp. 55-58) and §2.4
%                 (types of system calls, pp. 58-66), page-verified via
%                 master_supporting_docs/CS502/supporting_books/SilberschatzOSC/index.md;
%                 Bovet & Cesati, Understanding the Linux Kernel, 3rd ed. —
%                 Ch. 10 (system calls, pp. 398-419), page-verified via
%                 master_supporting_docs/CS502/supporting_books/LinuxKernel/index.md.
%
% Compile with: cd Notes/CS502 &&
%   TEXINPUTS=../../Preambles;$TEXINPUTS xelatex -interaction=nonstopmode 03-system-calls-notes.tex
%   BIBINPUTS=../..;$BIBINPUTS bibtex 03-system-calls-notes
%   TEXINPUTS=../../Preambles;$TEXINPUTS xelatex -interaction=nonstopmode 03-system-calls-notes.tex  (x2)
% =============================================================================

\documentclass[11pt]{article}
\usepackage[margin=1in]{geometry}
\usepackage{array}
\input{../../Preambles/header}
\coursecode{CS502}

% Chapter-style numbering: this is Week 3's chapter.
\renewcommand{\thesection}{3.\arabic{section}}
\renewcommand{\thefigure}{3.\arabic{figure}}
\renewcommand{\thetable}{3.\arabic{table}}

\newtheorem{example}{Example}
\renewcommand{\theexample}{3.\arabic{example}}

\newenvironment{definitionbox}[1]{%
  \par\noindent\textbf{Definition (#1).}\ \itshape
}{\par\vspace{0.3em}}

\title{System Calls\\\large Lecture Notes --- Week 3\\
  \normalsize Mechanics, Parameter Passing, and Types}
\author{Auqib Hamid Lone\\[0.3em]
  \emph{Department of Computer Science \& Engineering}, \emph{GCET Ganderbal}}
\date{\today}

\begin{document}
\maketitle

\section{The Command That Keeps Knocking}

Run the command

\begin{center}
\texttt{cp notes.txt backup.txt}
\end{center}

\noindent
and four things happen in order: the program opens one file, creates another,
copies every byte from the first into the second, and closes both. The command
looks like a single action at the shell prompt, but it is not a single action
inside the machine. The question that organizes this chapter is the one the
command keeps raising: \emph{how many times does the program cross from user
mode into the kernel, and what does it say each time?}

The answer cannot be found by reading the source code of \texttt{cp} alone.
\texttt{cp} is ordinary user code. It cannot touch the disk, modify the page
tables, or read another process's memory, because the hardware forbids those
operations in user mode. Every privileged operation it needs must be requested
through a \key{system call}. Each request has four parts: a name, a number, a
set of arguments, and a return value. The rest of the chapter is about those
four parts --- how they are expressed, how they cross the boundary, and how the
kernel decides whether to honour them.

We will keep \texttt{cp notes.txt backup.txt} as the running thread. Every
mechanism we introduce will be tied back to a step in that trace.

\subsection*{Roadmap}

Four questions frame the chapter, and they are also the four questions a
programmer must be able to answer about any operating system:

\begin{enumerate}
  \item \key{What is a system call?} It is the kernel's programming interface
        --- the boundary between a program and the services the OS provides.
  \item \key{How does control cross the boundary?} Through the trap and the
        system-call table.
  \item \key{How do arguments travel?} Through registers, a memory block, or
        the stack.
  \item \key{What kinds of calls exist?} The six standard categories:
        process control, file management, device management, information
        maintenance, communication, and protection.
\end{enumerate}

The division of labour is worth stating once. The \emph{mechanism} that carries
a request across the boundary is hardware; the \emph{interface} the programmer
uses is software; and the \emph{taxonomy} of calls is the vocabulary the rest of
the course will use.

\section{Two Views of the Same Service}

\noindent\muted{\textit{Reading: Silberschatz et al., Ch.~2 \S2.3, pp.~55--58.}}

Programmers rarely invoke a system call directly. They call a library function,
and the library eventually executes the trap. This is not an accident of
history; it is the normal layering of a C program, and it produces two views of
one service. Table~\ref{tab:two-views} sets the two views side by side.

\begin{table}[h]
\centering
\small
\renewcommand{\arraystretch}{1.3}
\begin{tabular}{@{}>{\raggedright\arraybackslash}p{0.42\linewidth}>{\raggedright\arraybackslash}p{0.42\linewidth}@{}}
\toprule
\textbf{Programmer's view (API)} & \textbf{Kernel's view (system call)} \\
\midrule
\texttt{fopen}, \texttt{fread}, \texttt{fwrite}, \texttt{fclose} &
\texttt{open}, \texttt{read}, \texttt{write}, \texttt{close} \\
Standard C library; buffered; portable &
Kernel entry points; unbuffered; architecture-specific \\
Returns a \texttt{FILE *} or a count &
Returns a file descriptor or $-1$ with \texttt{errno} \\
\bottomrule
\end{tabular}
\caption{Two views of the same file service. The API is what the programmer
writes; the system call is what the kernel executes.}
\label{tab:two-views}
\end{table}

The left column is a \emph{convenience layer}. It knows about buffering, it
hides the details of the machine's instruction set, and it offers a uniform
interface across operating systems. The right column is the \emph{actual
request} to the kernel: an entry point with a fixed number, a fixed calling
convention, and a return value the kernel can produce.

The two columns are not in one-to-one correspondence. One library call may
issue zero, one, or many system calls, because buffering changes the count.
A call that reads from a buffer the library already filled may never trap at
all; a call that fills a buffer may trap once for a large transfer; a call that
flushes may trap for each underlying write. This is why counting the system
calls in a program requires looking at the library's behaviour, not just at the
number of function calls in the source. Silberschatz et al.\ describe this
layering in Ch.~2 \S2.3 (pp.~55--58) \cite{SilberschatzOSC_operating_system_concepts}.

\section{What Exactly Is a System Call?}

\begin{definitionbox}{System call}
A \key{system call} is the programmatic request from a user program to a
kernel service. It is the boundary of the operating system's programming
interface.
\end{definitionbox}

The definition sounds like a function call, and at the source level it looks
like one: a name, parentheses, arguments, and a return value. But it is
\emph{not} an ordinary function call. Three differences matter.

First, it \emph{changes the privilege level}. The CPU leaves user mode, enters
kernel mode, runs kernel code, and returns to user mode. An ordinary function
call never changes the privilege level; caller and callee execute with the same
rights.

Second, the kernel executes the request \emph{on behalf of the calling process}
and \emph{with the kernel's privileges}. The process is still the owner of the
request --- the open file will belong to it, the bytes read will land in its
buffer --- but the code that performs the operation is trusted kernel code.

Third, the kernel \emph{validates the request before acting on it}. The user
program supplies arguments, but the kernel does not assume they are sensible.
This third difference is the one that turns the boundary from a convenience
into a security mechanism, and we will return to it repeatedly.

One warning about terminology: the same phrase ``system call'' is used for
three different things --- the API function the programmer calls, the trap
instruction that crosses the boundary, and the kernel routine that ultimately
runs. Context tells you which layer is meant. When this chapter says ``the
system call returns a file descriptor,'' it means the kernel routine; when it
says ``the system call is invoked with \texttt{syscall},'' it means the
instruction. The distinction is not pedantry: conflating the layers is the
source of most of the misconceptions at the end of the chapter.

\section{Why Not Link the Kernel Into Every Program?}

If a system call is just a request for a service, why not avoid the boundary
altogether? Why not link the kernel's code into every program, the way a
library is linked, and let each program call the routines directly? The answer
is that four properties would be destroyed at once.

\begin{itemize}
  \item \textbf{Protection.} User code would run with kernel privileges. A
        single buggy program could then corrupt the machine or overwrite the
        kernel itself.
  \item \textbf{Isolation.} One buggy program could overwrite another
        program's memory or the kernel's data structures, because nothing would
        prevent it from addressing them.
  \item \textbf{Portability.} The same API can sit on different kernels and
        different instruction sets only because the boundary hides the
        implementation. Linking the kernel in would bind every program to one
        kernel's internals.
  \item \textbf{Evolution.} The kernel can change its internal data structures
        and algorithms without recompiling every application, because the
        application sees only the stable interface, not the implementation.
\end{itemize}

The trade-off is stated just as plainly: every system call costs a mode switch.
The design goal is therefore to keep the boundary \emph{narrow} --- few, well
chosen entry points --- and the calls \emph{efficient}. The rest of the chapter
is, in a sense, the engineering of that trade-off: how to cross the boundary
with as little overhead as possible while still validating everything that
crosses.

\section{Why \texttt{printf} Is Not a System Call}

A useful test of the API/system-call distinction is the most familiar C
function of all. \texttt{printf} is not a system call. What does it eventually
call, and why does the C library buffer the output first?

\texttt{printf} formats its arguments into a \emph{user-space buffer}. The
formatting is ordinary library code running in user mode; none of it needs the
kernel. When the buffer fills, or when the program (or the library) decides to
flush it, the library calls \texttt{write}, and \texttt{write} traps into the
kernel. The kernel then copies the buffer to the destination.

Buffering exists to reduce the number of expensive mode switches. If every
\texttt{printf} trapped immediately, a program printing a line in twenty small
pieces would pay for twenty crossings. With a buffer, the same twenty pieces
become one \texttt{write} and one crossing. The bytes delivered are the same;
the number of boundary crossings is not. Table~\ref{tab:two-views}'s warning
that one library call may issue zero, one, or many system calls is exactly this
effect.

\clearpage
\section{The Trap Path}

\noindent\muted{\textit{Reading: Bovet \& Cesati, Ch.~10, pp.~398--419.}}

The trap instruction is the hinge of the whole mechanism: one instruction moves
the CPU from user mode to kernel mode and hands the kernel a number. Figure
\ref{fig:trap-path} shows the path that instruction starts.

\begin{figure}[h]
\centering
\begin{tikzpicture}[every node/.style={font=\small}]
  \node[dag-node, minimum width=2.8cm] (user) at (-5.2,0) {User code};
  \node[dag-node, minimum width=2.8cm] (lib) at (-1.8,0) {Library stub};
  \node[dag-node, minimum width=2.8cm] (trap) at (1.8,0) {Trap};
  \node[dag-node, minimum width=2.8cm] (table) at (5.2,0) {Syscall table};
  \node[dag-node, minimum width=2.8cm] (svc) at (5.2,-2.0) {Kernel handler};
  \draw[observed-edge] (user) -- (lib);
  \draw[observed-edge] (lib) -- (trap);
  \draw[observed-edge] (trap) -- (table);
  \draw[observed-edge] (table) -- (svc);
  \draw[observed-edge] (svc) -- ++(-5.2,0) -| (user);
\end{tikzpicture}
\caption{The trap path. User code calls a library stub; the stub executes the
trap; the kernel indexes its table by number, runs the handler, and returns to
the user program. The mechanism is described in Silberschatz et al., Ch.~2
\S2.3, pp.~55--58 \cite{SilberschatzOSC_operating_system_concepts}, and in
concrete Linux detail in Bovet \& Cesati, Ch.~10, pp.~398--419
\cite{BovetCesati_understanding_linux_kernel}.}
\label{fig:trap-path}
\end{figure}

The path has three stages. In the first stage, the library stub places the
\key{system-call number} and the arguments in registers. The number is what
identifies the requested service; the registers are where the kernel will look
for the arguments. In the second stage, the trap instruction saves the user
program counter (PC) and flags, switches the CPU to kernel mode, and jumps to
the kernel's trap entry point. Saving the PC is what makes a return possible:
the kernel must be able to resume the user program at the instruction after the
trap. In the third stage, the kernel indexes the \key{system-call table} by the
number it received and calls the corresponding handler. When the handler
finishes, the kernel restores the user context and returns.

The arrows in Figure~\ref{fig:trap-path} are drawn as a path, not as a cycle,
because the user program does not continue while the kernel works. Control has
moved to the kernel; the user code resumes only after the handler returns.

\section{Worked Example: \texttt{write(fd, buf, n)}}

The general path becomes concrete when we follow a single call. Consider

\begin{center}
\texttt{write(fd, buf, n)}
\end{center}

\noindent
where \texttt{fd} is a file descriptor, \texttt{buf} is the address of a
user buffer, and \texttt{n} is the number of bytes to write.

\begin{example}
\textup{\textbf{The seven steps of a \texttt{write} system call.}}
\begin{enumerate}
  \item \texttt{fd}, \texttt{buf}, and \texttt{n} are placed in the argument
        registers. This is still user-mode code; it is the library stub
        arranging the arguments where the kernel's calling convention expects
        them.
  \item The system-call number for \texttt{write} is placed in a designated
        register. The number and the arguments together describe the entire
        request: which service, and on what data.
  \item The stub executes \texttt{syscall} (x86-64) or \texttt{int 0x80}
        (legacy x86). This is the instruction that actually crosses the
        boundary.
  \item Hardware switches to kernel mode, saves the return address, and enters
        the trap handler. The return address is the user PC, saved so that the
        program can resume after the kernel finishes.
  \item The kernel validates \texttt{fd}, checks that \texttt{buf} lies in
        user memory, and performs the write. Validation comes first: the
        kernel will not copy to or from an address it has not checked.
  \item The return value is placed in a register; the kernel restores user mode
        and returns. The calling program continues at the saved return address.
  \item On error, the library sets \texttt{errno} and returns $-1$. The kernel
        signalled the error through its return value; the library translates it
        into the C convention the programmer expects.
\end{enumerate}
\end{example}

Steps 1--2 run in user mode, steps 3--6 run in the kernel, and step 7 runs back
in user mode. The boundary is crossed on entry (steps 3--4) and again on
return (step 6); the seven steps therefore contain one complete round trip. Bovet and
Cesati follow this path in detail for Linux in Ch.~10 (pp.~398--419)
\cite{BovetCesati_understanding_linux_kernel}, and Silberschatz et al.\ give
the architecture-neutral account in Ch.~2 \S2.3
\cite{SilberschatzOSC_operating_system_concepts}.

\section{The System-Call Table}

\begin{definitionbox}{System-call table}
A kernel array indexed by system-call number. Each entry is a function pointer
to the handler for that service.
\end{definitionbox}

The table is the second half of the hinge. The trap instruction carries only a
number; it does not carry a function name or an address chosen by the user
program. The table turns that number into the right kernel routine, and it does
so without a long search: the number is an index, so the kernel can go directly
to the entry.

Four properties follow from this design.

\begin{itemize}
  \item The number is the \key{ABI} (application binary interface): it must
        stay stable across kernel versions. A compiled program contains the
        number, not the name, so changing the number would break the program.
  \item Adding a call means adding a table entry and a handler --- not changing
        the trap instruction. The instruction is generic; the meaning lives in
        the table.
  \item The handler runs in kernel mode with the kernel's privileges, which is
        why the validation in step~5 of the worked example is possible at all.
  \item The table is small and indexed, so the cost of dispatch is a lookup
        rather than a search through names.
\end{itemize}

On Linux x86-64, \texttt{write} is system call 1, and the table lives in the
kernel's trap entry code. The number is not a historical accident to be
memorized for its own sake; it is the concrete form of the contract between a
compiled binary and the kernel it runs on.

\section{Return Values and \texttt{errno}}

A system call must report three kinds of outcome: success with a value, success
with no data (end of file), and failure. The kernel reports them through a
register. The value in that register is either a count, a file descriptor, or
an error code.

The convention is that a \emph{negative} return value signals an error. The
library wrapper translates that negative value into \texttt{-1} and sets
\texttt{errno} to a code identifying the failure. The caller must check the
return value, because the kernel cannot print a message for it.

The same discipline governs the arguments. The kernel never trusts user
pointers or lengths; it validates them before copying. A pointer that is
syntactically valid --- a number of the right width --- may still point
somewhere the process is not allowed to touch. Validation is the kernel's job,
and it happens before any data moves.

\begin{example}
\textup{\textbf{The three outcomes of \texttt{read}.}} A \texttt{read} call returns the
number of bytes read on success. It returns \texttt{0} at end of file, which is
a successful call that transferred no bytes. It returns $-1$ on error, with
\texttt{errno} set to a code such as \texttt{EFAULT} (a bad user address) or
\texttt{EBADF} (a bad file descriptor), or a similar value. The three cases are
distinguished by the return value, so a correct program must inspect it rather
than assume the call succeeded.
\end{example}

\section{Three Ways to Pass Arguments}

\noindent\muted{\textit{Reading: Silberschatz et al., Ch.~2 \S2.3, pp.~55--58.}}

A system call is a function call across a privilege boundary. The arguments
must cross too --- and the kernel must not trust any of them. There are three
standard mechanisms for getting the arguments across, and Figure
\ref{fig:three-ways} arranges them by their trade-off.

\begin{figure}[h]
\centering
\begin{tikzpicture}[every node/.style={font=\small}]
  \node[dag-node, minimum width=3.4cm, minimum height=1.1cm] (reg) at (-5.0,0)
    {Registers};
  \node[dag-node, minimum width=3.4cm, minimum height=1.1cm] (block) at (0,0)
    {Block / table\\in memory};
  \node[dag-node, minimum width=3.4cm, minimum height=1.1cm] (stack) at (5.0,0)
    {Stack};
  \node[text=primary-gold, align=center] at (-5.0,-1.25) {fast\\few arguments};
  \node[text=primary-gold, align=center] at (0,-1.25) {many arguments\\one pointer};
  \node[text=primary-gold, align=center] at (5.0,-1.25) {natural for\\C calling convention};
\end{tikzpicture}
\caption{The three standard mechanisms for passing system-call arguments.
Silberschatz et al.\ describe all three in Ch.~2 \S2.3, pp.~55--58
\cite{SilberschatzOSC_operating_system_concepts}.}
\label{fig:three-ways}
\end{figure}

Silberschatz describes all three; real ABIs choose one or a mix. The choice is
not cosmetic, because the kernel must know where to find the arguments before
it can validate them. A mechanism that is fast but limited, a mechanism that
scales to many arguments but requires an extra indirection, and a mechanism
that fits an existing calling convention are three different answers to the
same engineering question.

\section{Method 1: Registers}

The caller places the first few arguments in designated registers. This is the
fastest mechanism: the values are already in the CPU, so no memory access is
needed beyond the registers that were already saved on entry. The limitation is
equally clear: the method is bounded by the number of argument registers. A
call with more arguments than the ABI reserves cannot use this method alone.

\begin{example}
\textup{\textbf{Linux x86-64 \texttt{read(fd, buf, count)}.}} The convention for this
call is:
\begin{center}
\begin{tabular}{@{}ll@{}}
\texttt{rax} & system-call number for \texttt{read} \\
\texttt{rdi} & \texttt{fd} \\
\texttt{rsi} & \texttt{buf} (user virtual address) \\
\texttt{rdx} & \texttt{count} \\
\end{tabular}
\end{center}
The number travels in \texttt{rax}; the three arguments travel in \texttt{rdi},
\texttt{rsi}, and \texttt{rdx}. The choice of registers is part of the ABI, so
the library stub and the kernel agree on it without negotiation.
\end{example}

Register passing does not relax validation. The kernel still checks that
\texttt{buf} is a legal user address; putting the address in \texttt{rsi} says
where the buffer is, not that the program is entitled to use it.

\section{Method 2: Block or Table in Memory}

When the argument list is too large for the reserved registers, the caller can
store the arguments in a memory block and pass the \emph{address} of the block
in a register. The kernel copies the block into kernel space and validates
every field. One register now carries an arbitrary number of arguments, at the
price of one indirection and a copy.

\begin{example}
\textbf{A call with many arguments.} \texttt{open(path, flags, mode)} can be
passed as a table holding the pointer to \texttt{path}, the value of
\texttt{flags}, and the value of \texttt{mode}. A single register holds the
address of the table; the kernel reads the table once, copies it into kernel
memory, and validates the pointer to \texttt{path} along with the other fields.
\end{example}

The pointer to the block is itself untrusted. The kernel must verify that the
\emph{whole} block lies in the caller's address space before it reads any field,
because a block that starts in user memory and continues into kernel memory
would be just as dangerous as a direct pointer into the kernel.

\section{Method 3: Stack}

The caller can also push the arguments onto the user stack, from which the
kernel reads them using the saved user stack pointer. This is natural for
languages and calling conventions that already pass arguments on the stack: the
stub does not have to rearrange anything, because the arguments are already
laid out in the expected place.

Whichever method is used, the kernel must satisfy the same three requirements:

\begin{enumerate}
  \item \emph{Locate} the arguments --- know which registers, block, or stack
        slots hold them.
  \item \emph{Copy} them into kernel memory, so that the kernel is not reading
        a user buffer that the user program could change underneath it.
  \item \emph{Validate} every pointer and length before using it, so that a
        malformed argument produces an error rather than a corrupted kernel.
\end{enumerate}

These three steps are the invariant behind the three mechanisms. The mechanisms
differ in where the arguments start; they do not differ in what the kernel must
do once it has found them.

\section{Pointer Validation: Never Trust User Memory}

A user pointer is just a number, and it may point anywhere: into kernel memory,
into another process's memory, or into unmapped memory that will fault on
access. The kernel therefore checks the range against the process's address
space before it uses the pointer. It then copies data with special routines
(\texttt{copy\_from\_user}/\texttt{copy\_to\_user} on Linux) that handle
faults safely. A bad pointer returns \texttt{EFAULT}; it does not crash the
kernel.

The security rule is absolute: every argument is untrusted input. Validation
happens in kernel mode, before the kernel acts on the request. This is the
reason the boundary exists. If the kernel trusted a user pointer, a program
could pass the address of a kernel data structure to \texttt{read} and receive
the kernel's memory as output, or pass it to \texttt{write} and overwrite kernel
state. Dereferencing the pointer directly would make the protection of the
previous section meaningless.

A concrete test makes the rule vivid. Suppose a user program passes a pointer
to a kernel address to \texttt{read}. The kernel must not follow the pointer.
It checks that the range lies in the user address space; if it does not, the
call fails with \texttt{EFAULT}. The program learns that its request was
invalid, and the kernel remains intact. The check is not an optimization; it is
the boundary being enforced.

\section{The Six Categories}

\noindent\muted{\textit{Reading: Silberschatz et al., Ch.~2 \S2.4, pp.~58--66.}}

System calls are not a random list. They fall into six categories, and every
command you run is a combination of them. Table~\ref{tab:six-categories} gives
the taxonomy.

\begin{table}[h]
\centering
\small
\renewcommand{\arraystretch}{1.25}
\begin{tabular}{@{}>{\raggedright\arraybackslash}p{0.24\linewidth}>{\raggedright\arraybackslash}p{0.35\linewidth}>{\raggedright\arraybackslash}p{0.31\linewidth}@{}}
\toprule
\textbf{Category} & \textbf{Purpose} & \textbf{Typical calls} \\
\midrule
Process control & create, terminate, wait, load &
\texttt{fork}, \texttt{exec}, \texttt{wait}, \texttt{exit} \\
File management & create, open, read, write, close &
\texttt{open}, \texttt{read}, \texttt{write}, \texttt{close}, \texttt{lseek} \\
Device management & request/release, read/write, control &
\texttt{ioctl}, \texttt{read}, \texttt{write} \\
Information maintenance & get/set time, process, file attributes &
\texttt{getpid}, \texttt{alarm}, \texttt{sleep}, \texttt{stat} \\
Communication & create/delete connections, send/receive &
\texttt{pipe}, \texttt{mmap}, \texttt{socket}, \texttt{send}, \texttt{recv} \\
Protection & set/get permissions &
\texttt{chmod}, \texttt{umask}, \texttt{chown}, \texttt{setuid} \\
\bottomrule
\end{tabular}
\caption{The six categories of system calls. Silberschatz et al., Ch.~2
\S2.4, pp.~58--66 \cite{SilberschatzOSC_operating_system_concepts}.}
\label{tab:six-categories}
\end{table}

The categories are not a hierarchy and they are not mutually exclusive in
practice: a single operation may touch several of them. They are a vocabulary
for describing what a call is \emph{for}, and the rest of the course will use
that vocabulary constantly.

\section{Process Control}

The process-control calls manage the life cycle of a process.

\begin{itemize}
  \item \texttt{fork}: create a child process that is a copy of the parent.
        The child inherits the parent's open-file table.
  \item \texttt{exec}: replace the current process image with a new program,
        so the process runs different code.
  \item \texttt{wait}: block until a child terminates and collect its status.
        This is how a parent learns that a child has finished and whether it
        succeeded.
  \item \texttt{exit}: terminate the calling process and release its resources.
\end{itemize}

The shell's job is the canonical combination. A shell \texttt{fork}s, the child
\texttt{exec}s the command, and the parent \texttt{wait}s. When you type
\texttt{cp notes.txt backup.txt}, the shell creates a child, the child replaces
itself with the \texttt{cp} program through \texttt{exec}, and \texttt{cp}
itself is the child after \texttt{exec}. The shell is still waiting; that is
why the prompt returns only when the copy finishes.

\section{File Management}

The file-management calls operate on files and on the descriptors that
represent them.

\begin{itemize}
  \item \texttt{open}: create or open a file, returning a file descriptor.
  \item \texttt{read}/\texttt{write}: transfer bytes between user memory and
        the file.
  \item \texttt{lseek}: move the file offset without transferring data.
  \item \texttt{close}: release the descriptor and flush state.
  \item \texttt{stat}: read metadata (size, owner, permissions, times).
\end{itemize}

\begin{definitionbox}{File descriptor}
A small integer index into the process's per-process open-file table --- not
the file itself.
\end{definitionbox}

The distinction in the definition matters. The descriptor is a handle, and it
is meaningful only in the process that owns it; the open-file table entry
behind the handle is what actually refers to the file.

\section{Device, Information, Communication, Protection}

The remaining four categories cover everything that is not process or file
management.

\paragraph{Device management.} \texttt{ioctl} controls device-specific
behaviour; \texttt{read} and \texttt{write} move data. Devices are often
accessed through the same descriptor interface as files, which is why the same
\texttt{read} and \texttt{write} calls appear in both the file-management and
device-management rows of Table~\ref{tab:six-categories}. The uniformity is a
deliberate design choice: a program that knows how to read a file already knows
how to read from many devices.

\paragraph{Information maintenance.} \texttt{getpid}, \texttt{alarm},
\texttt{sleep}, and \texttt{time} retrieve or manipulate information about the
process and the system; \texttt{stat} retrieves file metadata. These calls do
not move bulk data; they query or set small pieces of state.

\paragraph{Communication.} \texttt{pipe}, \texttt{mmap}, \texttt{socket},
\texttt{send}, and \texttt{recv} move data between processes or across a
network.

\paragraph{Protection.} \texttt{chmod}, \texttt{umask}, \texttt{chown}, and
\texttt{setuid} set or query access rights. They do not read or write file
contents; they change who is allowed to do so.

\section{Worked Example: \texttt{cp} Across the Categories}

The six categories stop being abstract when we lay the running thread out
against them.

\begin{example}
\textup{\textbf{\texttt{cp notes.txt backup.txt} as a sequence of system calls.}}
\begin{enumerate}
  \item \textbf{Process control:} the shell \texttt{fork}s and the child
        \texttt{exec}s \texttt{cp}.
  \item \textbf{File management:} \texttt{open} is called on
        \texttt{notes.txt} with \texttt{O\_RDONLY}, and on
        \texttt{backup.txt} with \texttt{O\_WRONLY|O\_CREAT}, which creates or
        truncates the destination.
  \item \textbf{File management:} a loop of \texttt{read} into a buffer and
        \texttt{write} to the destination moves the bytes.
  \item \textbf{File management:} \texttt{close} releases both descriptors.
  \item \textbf{Process control:} \texttt{exit(0)} ends the child; the shell's
        \texttt{wait} collects the status.
\end{enumerate}
\end{example}

One shell command, ten or more system calls. The count is worth deriving
explicitly, because it is the answer to the question the chapter opened with.
There is one \texttt{fork}, one \texttt{exec}, two \texttt{open} calls, two
\texttt{close} calls, one \texttt{exit}, and one \texttt{wait} --- eight calls
before any data moves. At least one \texttt{read} and one \texttt{write} are
needed even for a one-byte file, which brings the minimum to ten. A real file
needs many more iterations of the read/write loop, so the count grows with the
file size. The categories are not abstract: they are the steps in the trace.

\section{Why \texttt{fork} Is Process Control}

A final classification check. Which category does \texttt{fork} belong to, and
why is it not file management even though the child may open files?

\texttt{fork} creates a new process, so it is process control. The child
inherits the parent's open-file table, but \texttt{fork} itself does not create
or access a file. The category describes the \emph{operation}, not the
resources the process may use afterward. This is the rule to apply to every
call in the taxonomy: classify the operation, not its consequences.

\section{Numbers, ABI, and Compatibility}

\noindent\muted{\textit{Reading: Bovet \& Cesati, Ch.~10, pp.~398--419;
Silberschatz et al., Ch.~2 \S2.3, pp.~55--58.}}

The interface is stable, but every call has a cost. Good systems programmers
batch calls and know when the boundary is being crossed. Before turning to the
cost, it is worth seeing why the interface is stable in the first place.

Each architecture defines a system-call number for each service. The numbers
are part of the \key{application binary interface}. A kernel cannot renumber
existing calls without breaking every compiled program, because the number is
embedded in the binary, not looked up by name at run time. New calls are
therefore added at the end of the table, and old calls remain for
compatibility.

\begin{definitionbox}{Application binary interface}
The contract between compiled code and the kernel: the calling convention and,
for system calls, the number assigned to each service.
\end{definitionbox}

The consequence is that the system-call table is a long-lived contract between
the kernel and every binary that runs on it. The contract can grow, but it
cannot be silently rewritten.

\section{The Cost of Crossing the Boundary}

A system call is not free. The cost has four components:

\begin{itemize}
  \item Save the user context and switch to kernel mode.
  \item Index the table and enter the handler.
  \item Validate the arguments and copy data across the boundary.
  \item Perform the operation, then restore the user context.
\end{itemize}

Each component is small, but they add up, and the copy in the third component
is proportional to the amount of data moved. This is why batching matters.
Copying 1\,MB in one \texttt{read} costs one mode switch; copying it in 1-byte
chunks costs roughly a million. The bytes are the same; the crossings are not.
The same arithmetic applies to writes, and it is the reason library buffering
(recall \texttt{printf}) is worth the complexity.

\texttt{strace} shows the actual calls a program makes. It is the practical way
to see the boundary in action: instead of guessing how many times a program
traps, you can read the sequence of calls it really issued.

\section{Common Misconceptions}

Table~\ref{tab:misconceptions} collects the four errors that the chapter has
been building the tools to correct.

\begin{table}[h]
\centering
\small
\renewcommand{\arraystretch}{1.25}
\begin{tabular}{@{}>{\raggedright\arraybackslash}p{0.44\linewidth}>{\raggedright\arraybackslash}p{0.44\linewidth}@{}}
\toprule
\textbf{Wrong} & \textbf{Right} \\
\midrule
``\texttt{printf} is a system call.'' &
\texttt{printf} uses buffered \texttt{write} calls. \\
``A system call is just a function call.'' &
It changes privilege level and uses a kernel entry point. \\
``The kernel can trust the user's pointer.'' &
User pointers and lengths are untrusted and validated. \\
``Every library call traps into the kernel.'' &
Many library calls stay entirely in user space. \\
\bottomrule
\end{tabular}
\caption{Four common misconceptions about system calls, and the correction for
each.}
\label{tab:misconceptions}
\end{table}

Each correction is a one-line summary of a mechanism developed earlier:
the first of the API/system-call distinction (\S3.2 and \S3.5), the second of
the trap path (\S3.6), the third of pointer validation (\S3.14), and the fourth
of the fact that one library call may issue zero, one, or many system calls
(\S3.2).

\section{Summary}

A system call is the program's request to the kernel; it is invoked through a
trap, not an ordinary call. The system-call number selects a handler from the
kernel's dispatch table, and the table is a long-lived ABI contract. Arguments
travel in registers, in a memory block, or on the stack, and every pointer and
length is validated before the kernel acts. The six categories --- process
control, file management, device management, information maintenance,
communication, and protection --- organize the interface, and a single command
such as \texttt{cp notes.txt backup.txt} is a combination of them. Mode
switches are expensive, so work is batched to reduce them; the library API sits
on top of the system calls and decides, through buffering, how many crossings a
program actually pays for.

The chapter's opening question now has a concrete answer. \texttt{cp} crosses
the boundary at least ten times for a one-byte file --- once per system call,
with a matching return crossing for each --- and the number grows with the file
size because the read/write loop repeats. What it says each time is the
system-call number, followed by arguments the kernel validates before it acts.

\section{Looking Ahead and Further Reading}

This chapter covered how a program names a kernel service, how the arguments
cross the boundary, and what kinds of calls exist. The next lecture moves up
one level: it examines the user-level programs built on top of these calls, and
the internal structures --- monolithic, layered, microkernel, modular, and
virtual machine --- that organize the kernel itself. The reading for that
lecture is Silberschatz et al., Ch.~2 \S2.5--2.8
\cite{SilberschatzOSC_operating_system_concepts}, and Stallings, Ch.~2 and
Ch.~4 \S4.3 \cite{StallingsOS_internals_design_principles}.

For this chapter, the page-cited sources are:

\begin{itemize}
  \item A.~Silberschatz, P.~B.~Galvin, and G.~Gagne, \emph{Operating System
        Concepts}, 8th ed. --- Ch.~2 \S2.3 (system calls, pp.~55--58) and
        \S2.4 (types of system calls, pp.~58--66).
  \item D.~P.~Bovet and M.~Cesati, \emph{Understanding the Linux Kernel}, 3rd
        ed. --- Ch.~10 (system calls, pp.~398--419), for the concrete
        x86/Linux entry and exit path.
  \item W.~Stallings, \emph{Operating Systems: Internals and Design
        Principles}, 6th ed. --- Ch.~3 (process description and control), for
        the process context in which calls execute.
\end{itemize}

\noindent
All three citations are page-verified through each book's \texttt{index.md} in
the course's supporting-books directory:
\begin{center}
\texttt{master\_supporting\_docs/CS502/supporting\_books/}
\end{center}

\section*{Self-Check}
\addcontentsline{toc}{section}{Self-Check}

\subsection*{Exercises}

\begin{enumerate}
  \item List the system calls in the \texttt{cp notes.txt backup.txt} trace,
        classify each by category, and state the minimum number of user-to-kernel
        crossings for a one-byte file. Explain why the count grows with the file
        size.
  \item A program calls \texttt{printf} one hundred times, one byte per call,
        and then exits. Does that imply one hundred \texttt{write} system calls?
        Explain using the user-space buffer and the conditions under which a
        flush occurs.
  \item A hypothetical system call takes seven arguments. Which of the three
        argument-passing mechanisms would you choose, and what must the kernel
        do before it uses the arguments? Be specific about the pointer.
  \item A user program passes a pointer to a kernel address to \texttt{read}.
        What should happen, and why can the kernel not simply dereference the
        pointer?
  \item A program copies a 4\,MB file in 4\,KB chunks, with one \texttt{read}
        and one \texttt{write} per chunk. Compare the number of system calls
        with a single 4\,MB \texttt{read} followed by a single 4\,MB
        \texttt{write}. Which cost dominates in each case?
\end{enumerate}

\subsection*{Solutions}

\begingroup\small
\begin{enumerate}
  \item The calls are \texttt{fork} (process control), \texttt{exec} (process
        control), two \texttt{open} calls (file management), at least one
        \texttt{read} and one \texttt{write} (file management), two
        \texttt{close} calls (file management), \texttt{exit} (process
        control), and the shell's \texttt{wait} (process control). For a
        one-byte file that is $1+1+2+1+1+2+1+1 = 10$ calls, hence at least ten
        user-to-kernel crossings (and ten matching kernel-to-user returns). A
        larger file repeats the read/write pair, so the count grows with the
        number of chunks, not with the number of commands.

  \item No. \texttt{printf} formats each byte into a user-space buffer. The
        library calls \texttt{write} only when the buffer fills or when it is
        flushed. If the buffer does not fill before exit, the hundred calls may
        produce far fewer than one hundred \texttt{write} calls --- possibly a
        single flush at exit. The output bytes are the same; the number of
        boundary crossings is not.

  \item Use a block or table in memory. The caller places the seven arguments
        in a block and passes the block's address in one register. The kernel
        locates the block, copies it into kernel memory, and validates every
        field --- in particular, it must verify that the whole block, and not
        just its first byte, lies in the caller's address space. Register
        passing alone would be limited by the number of argument registers;
        the stack is a viable alternative if the calling convention already
        places arguments there, but the block is the direct answer to ``too
        many arguments for the registers.''

  \item The call should fail with \texttt{EFAULT}. Dereferencing the pointer
        would let user code read or overwrite kernel data, which is exactly
        what the privilege boundary exists to prevent. The kernel checks that
        the range lies in the user address space before copying; if it does
        not, the request is rejected without touching the address.

  \item In 4\,KB chunks, $4\,\text{MB} = 4096\,\text{KB}$, so the program issues
        $4096/4 = 1024$ \texttt{read} calls and 1024 \texttt{write} calls:
        2048 system calls, hence 2048 mode-switch round trips. With one large
        \texttt{read} and one large \texttt{write}, it issues 2 system calls
        and 2 round trips. The chunked version pays the fixed cost of crossing
        the boundary 1024 times more often; the large-transfer version pays it
        twice, but must hold 4\,MB in memory at once. The fixed cost of the
        crossing dominates the chunked case; the buffer size dominates the
        large-transfer case.
\end{enumerate}
\endgroup

\begin{thebibliography}{9}
\bibitem{SilberschatzOSC_operating_system_concepts}
A.~Silberschatz, P.~B.~Galvin, and G.~Gagne, \emph{Operating System Concepts},
8th ed., John Wiley and Sons --- Ch.~2 \S2.3 (system calls, pp.~55--58) and
\S2.4 (types of system calls, pp.~58--66). \emph{(Page-cited; index built
2026-10-14.)}
\bibitem{StallingsOS_internals_design_principles}
W.~Stallings, \emph{Operating Systems: Internals and Design Principles}, 6th
ed., Prentice Hall --- Ch.~2, Ch.~3 (process description and control), and
Ch.~4 \S4.3. \emph{(Page-cited; index built 2026-10-14.)}
\bibitem{BovetCesati_understanding_linux_kernel}
D.~P.~Bovet and M.~Cesati, \emph{Understanding the Linux Kernel}, 3rd ed.,
O'Reilly Media, 2006 --- Ch.~10 (system calls, pp.~398--419), for the concrete
x86/Linux entry and exit path. \emph{(Page-cited; index built 2026-10-14.)}
\end{thebibliography}

\end{document}
